% GENERATED FILE. Edit canonical lessons, then run npm run generate:books.
% canonical-source-hash: fnv1a64:35bf3b6086b3c6b9
% canonical-lessons: PT-C272-durar, PT-C272-prolongar, PT-C272-instante, PT-C272-periodo, PT-C272-imediatamente

\chapter{Quanto tempo durou}
\label{ch:pt-quanto-tempo-durou}
\hlchaptermodality{272}

\begin{chapteropening}
\textbf{By the end of this chapter:} \emph{I can say how long something lasted, that it was prolonged, and that something happened in an instant.}

Complete the last lesson of chapter 272: I can say how long something lasted, that it was prolonged, and that something happened in an instant.
\end{chapteropening}

\section[durar]{durar — to last}

\label{lesson:PT-C272-durar}



\begin{quote}
Before the new one: say the Portuguese for the event, then the Portuguese for the incident.
\end{quote}

\subsection*{You'll want to know: durar}
\textbf{durar} — \textquotedblleft{}to last\textquotedblright{}.

\begin{quote}
A viagem durou cinco horas.
\end{quote}

\begin{quote}
\emph{The journey lasted five hours.}
\end{quote}

\emph{Durar muito} — to last a long time; \emph{durar pouco} — not to last long.

\begin{cousinweb}
\emph{durar} is inherited from Latin \emph{dūrāre} \textquotedblleft{}to harden, to hold out, to last\textquotedblright{} (Lewis \& Short), built on \emph{dūrus} \textquotedblleft{}hard\textquotedblright{}.
\end{cousinweb}

\begin{grammarlens}[title={What you've built}]
One more word for telling, explaining and describing at B1.
\end{grammarlens}

\subsection*{Your turn}
\begin{itemize}
\raggedright
  \item \emph{Say it:} \emph{durar}
  \item \emph{Say it:} \emph{durar}, once more
  \item \emph{Say it:} a sentence of your own with \emph{durar}
  \item \emph{Recall:} say the Portuguese for the event, then the Portuguese for the incident, then say \emph{durar} again
\end{itemize}

\begin{tcolorbox}[breakable,colback=teal!4,colframe=teal!35!black,title={Before you move on}]
What does durar mean? (\textquotedblleft{}To last\textquotedblright{}.) Say it once more.
\end{tcolorbox}
\section[prolongar]{prolongar — to prolong, to extend}

\label{lesson:PT-C272-prolongar}



\begin{quote}
Before the new one: say the Portuguese for the incident, then the Portuguese for to last.
\end{quote}

\subsection*{You'll want to know: prolongar}
\textbf{prolongar} — \textquotedblleft{}to prolong, to extend\textquotedblright{}.

\begin{quote}
Decidimos prolongar as férias mais uma semana.
\end{quote}

\begin{quote}
\emph{We decided to extend the holiday by another week.}
\end{quote}

For time and for length: \emph{prolongar uma estrada}.

\begin{cousinweb}
\emph{prolongar} was borrowed from Latin \emph{prōlongāre} \textquotedblleft{}to lengthen\textquotedblright{} (EtymDB), built on \emph{longus} \textquotedblleft{}long\textquotedblright{}.
\end{cousinweb}

\begin{grammarlens}[title={What you've built}]
One more word for telling, explaining and describing at B1.
\end{grammarlens}

\subsection*{Your turn}
\begin{itemize}
\raggedright
  \item \emph{Say it:} \emph{prolongar}
  \item \emph{Say it:} \emph{prolongar}, once more
  \item \emph{Say it:} a sentence of your own with \emph{prolongar}
  \item \emph{Recall:} say the Portuguese for the incident, then the Portuguese for to last, then say \emph{prolongar} again
\end{itemize}

\begin{tcolorbox}[breakable,colback=teal!4,colframe=teal!35!black,title={Before you move on}]
What does prolongar mean? (\textquotedblleft{}To prolong, to extend\textquotedblright{}.) Say it once more.
\end{tcolorbox}
\section[o instante]{o instante — the instant, the moment}

\label{lesson:PT-C272-instante}



\begin{quote}
Before the new one: say the Portuguese for to last, then the Portuguese for to prolong.
\end{quote}

\subsection*{You'll want to know: o instante}
\textbf{o instante} — \textquotedblleft{}the instant, the moment\textquotedblright{}.

\begin{quote}
Num instante, o céu ficou escuro.
\end{quote}

\begin{quote}
\emph{In an instant the sky went dark.}
\end{quote}

\emph{Um instante!} — one moment!

\begin{cousinweb}
\emph{instante} is a learned borrowing from Latin \emph{īnstāns} \textquotedblleft{}pressing, present\textquotedblright{} (EtymDB), from \emph{īnstāre} \textquotedblleft{}to stand upon, to press on\textquotedblright{}, built on \emph{stāre} \textquotedblleft{}to stand\textquotedblright{}.
\end{cousinweb}

\begin{grammarlens}[title={What you've built}]
One more word for telling, explaining and describing at B1.
\end{grammarlens}

\subsection*{Your turn}
\begin{itemize}
\raggedright
  \item \emph{Say it:} \emph{o instante}
  \item \emph{Say it:} \emph{o instante}, once more
  \item \emph{Say it:} a sentence of your own with \emph{o instante}
  \item \emph{Recall:} say the Portuguese for to last, then the Portuguese for to prolong, then say \emph{o instante} again
\end{itemize}

\begin{tcolorbox}[breakable,colback=teal!4,colframe=teal!35!black,title={Before you move on}]
What does o instante mean? (\textquotedblleft{}The instant, the moment\textquotedblright{}.) Say it once more.
\end{tcolorbox}
\section[o período]{o período — the period, the stretch of time}

\label{lesson:PT-C272-periodo}



\begin{quote}
Before the new one: say the Portuguese for to prolong, then the Portuguese for the instant.
\end{quote}

\subsection*{You'll want to know: o período}
\textbf{o período} — \textquotedblleft{}the period, the stretch of time\textquotedblright{}.

\begin{quote}
Vivi numa aldeia durante um período curto.
\end{quote}

\begin{quote}
\emph{I lived in a village for a short period.}
\end{quote}

At school, a term can also be \emph{um período}.

\begin{cousinweb}
\emph{período} is a learned borrowing from Latin \emph{periodus}, from Greek \emph{períodos} \textquotedblleft{}a going round, a cycle\textquotedblright{} (Lewis \& Short).
\end{cousinweb}

\begin{grammarlens}[title={What you've built}]
One more word for telling, explaining and describing at B1.
\end{grammarlens}

\subsection*{Your turn}
\begin{itemize}
\raggedright
  \item \emph{Say it:} \emph{o período}
  \item \emph{Say it:} \emph{o período}, once more
  \item \emph{Say it:} a sentence of your own with \emph{o período}
  \item \emph{Recall:} say the Portuguese for to prolong, then the Portuguese for the instant, then say \emph{o período} again
\end{itemize}

\begin{tcolorbox}[breakable,colback=teal!4,colframe=teal!35!black,title={Before you move on}]
What does o período mean? (\textquotedblleft{}The period, the stretch of time\textquotedblright{}.) Say it once more.
\end{tcolorbox}
\section[imediatamente]{imediatamente — immediately, at once}

\label{lesson:PT-C272-imediatamente}



\begin{quote}
Before the new one: say the Portuguese for the instant, then the Portuguese for the period.
\end{quote}

\subsection*{You'll want to know: imediatamente}
\textbf{imediatamente} — \textquotedblleft{}immediately, at once\textquotedblright{}.

\begin{quote}
O médico chegou imediatamente.
\end{quote}

\begin{quote}
\emph{The doctor arrived immediately.}
\end{quote}

\emph{Já} is the short, everyday way to say it.

\begin{cousinweb}
\emph{imediatamente} is built on \emph{imediato}, a learned word from Late Latin \emph{immediātus} \textquotedblleft{}with nothing in between\textquotedblright{}, from \emph{in-} \textquotedblleft{}not\textquotedblright{} and \emph{medius} \textquotedblleft{}middle\textquotedblright{} (Lewis \& Short).
\end{cousinweb}

\begin{grammarlens}[title={What you've built}]
One more word for telling, explaining and describing at B1.
\end{grammarlens}

\subsection*{Your turn}
\begin{itemize}
\raggedright
  \item \emph{Say it:} \emph{imediatamente}
  \item \emph{Say it:} \emph{imediatamente}, once more
  \item \emph{Say it:} a sentence of your own with \emph{imediatamente}
  \item \emph{Recall:} say the Portuguese for the instant, then the Portuguese for the period, then say \emph{imediatamente} again
\end{itemize}

\begin{tcolorbox}[breakable,colback=teal!4,colframe=teal!35!black,title={Before you move on}]
What does imediatamente mean? (\textquotedblleft{}Immediately, at once\textquotedblright{}.) Say it once more.
\end{tcolorbox}
